% GENERATED FILE. Edit canonical lessons, then run npm run generate:books.
% canonical-source-hash: fnv1a64:3a1ead8fb2dd5a19
% canonical-lessons: TA-C521-ilupparai, TA-C521-virippu, TA-C521-ottatai, TA-C521-talppal, TA-C521-puntotti, TA-R521-review-514-521

\chapter{Around the House}
\label{ch:ta-around-the-house}
\hlchaptermodality{521}

\begin{chapteropening}
\textbf{By the end of this chapter:} \emph{I can name more small things around the house.}

Complete the last lesson of chapter 521: I can name more small things around the house.
\end{chapteropening}

\section[\texorpdfstring{\ta{இழுப்பறை} (iḻuppaṟai)}{இழுப்பறை (iḻuppaṟai)}]{\ta{இழுப்பறை} (iḻuppaṟai) — a drawer}

\label{lesson:TA-C521-ilupparai}



\begin{quote}
Before the new one: say the Tamil for an electrical appliance, then the Tamil for to iron.
\end{quote}

\subsection*{You'll want to know: \ta{இழுப்பறை}}
\textbf{\ta{இழுப்பறை}} — \emph{iḻuppaṟai} — \textquotedblleft{}a drawer\textquotedblright{}.

\begin{quote}
\ta{சாவி} \ta{மேசையின்} \ta{இழுப்பறையில்} \ta{இருக்கிறது}.
\end{quote}

\begin{quote}
\emph{cāvi mēcaiyiṉ iḻuppaṟaiyil irukkiṟatu.}
\end{quote}

\begin{quote}
The key is in the desk drawer.
\end{quote}

Word for word \textquotedblleft{}a pull-room\textquotedblright{}: \ta{இழு} (\emph{iḻu}) \textquotedblleft{}to pull\textquotedblright{} + \ta{அறை} (\emph{aṟai}) \textquotedblleft{}a room\textquotedblright{}.

\begin{cousinweb}
\textbf{\ta{இழுப்பறை}} (\emph{iḻuppaṟai}) is a native compound: \ta{இழு} (\emph{iḻu}), \textquotedblleft{}to pull, to draw\textquotedblright{} (DEDR 504), and \ta{அறை} (\emph{aṟai}), \textquotedblleft{}a room, a compartment\textquotedblright{}.
\end{cousinweb}

\begin{grammarlens}[title={What you've built}]
One more word for describing what you lived and felt.
\end{grammarlens}

\subsection*{Your turn}
\begin{itemize}
\raggedright
  \item \emph{Say it:} \emph{iḻuppaṟai}
  \item \emph{Say it:} \emph{iḻuppaṟai}, once more
  \item \emph{Say it:} \emph{cāvi mēcaiyiṉ iḻuppaṟaiyil irukkiṟatu.} — The key is in the desk drawer.
  \item \emph{Recall:} say the Tamil for an electrical appliance, then the Tamil for to iron, then say \emph{iḻuppaṟai} again
\end{itemize}

\begin{tcolorbox}[breakable,colback=teal!4,colframe=teal!35!black,title={Before you move on}]
What does \ta{இழுப்பறை} mean? (\textquotedblleft{}A drawer\textquotedblright{}.) Say it once more.
\end{tcolorbox}
\section[\texorpdfstring{\ta{விரிப்பு} (virippu)}{விரிப்பு (virippu)}]{\ta{விரிப்பு} (virippu) — a bedspread; a sheet}

\label{lesson:TA-C521-virippu}



\begin{quote}
Before the new one: say the Tamil for to iron, then the Tamil for a drawer.
\end{quote}

\subsection*{You'll want to know: \ta{விரிப்பு}}
\textbf{\ta{விரிப்பு}} — \emph{virippu} — \textquotedblleft{}a bedspread; a sheet\textquotedblright{}.

\begin{quote}
\ta{அம்மா} \ta{கட்டிலில்} \ta{ஒரு} \ta{புதிய} \ta{விரிப்பை} \ta{விரித்தாள்}.
\end{quote}

\begin{quote}
\emph{ammā kaṭṭilil oru putiya virippai virittāḷ.}
\end{quote}

\begin{quote}
Mother spread a new sheet on the bed.
\end{quote}

It comes from \ta{விரி}, to spread out; the noun and the verb often go together, as here.

\begin{cousinweb}
\textbf{\ta{விரிப்பு}} (\emph{virippu}) is native Dravidian, a noun made from \ta{விரி} (\emph{viri}), \textquotedblleft{}to spread out, to unfold\textquotedblright{} (DEDR 5411), with the ending -\ta{ப்பு}.
\end{cousinweb}

\begin{grammarlens}[title={What you've built}]
One more word for describing what you lived and felt.
\end{grammarlens}

\subsection*{Your turn}
\begin{itemize}
\raggedright
  \item \emph{Say it:} \emph{virippu}
  \item \emph{Say it:} \emph{virippu}, once more
  \item \emph{Say it:} \emph{ammā kaṭṭilil oru putiya virippai virittāḷ.} — Mother spread a new sheet on the bed.
  \item \emph{Recall:} say the Tamil for to iron, then the Tamil for a drawer, then say \emph{virippu} again
\end{itemize}

\begin{tcolorbox}[breakable,colback=teal!4,colframe=teal!35!black,title={Before you move on}]
What does \ta{விரிப்பு} mean? (\textquotedblleft{}A bedspread; a sheet\textquotedblright{}.) Say it once more.
\end{tcolorbox}
\section[\texorpdfstring{\ta{ஒட்டடை} (oṭṭaṭai)}{ஒட்டடை (oṭṭaṭai)}]{\ta{ஒட்டடை} (oṭṭaṭai) — a cobweb}

\label{lesson:TA-C521-ottatai}



\begin{quote}
Before the new one: say the Tamil for a drawer, then the Tamil for a bedspread.
\end{quote}

\subsection*{You'll want to know: \ta{ஒட்டடை}}
\textbf{\ta{ஒட்டடை}} — \emph{oṭṭaṭai} — \textquotedblleft{}a cobweb\textquotedblright{}.

\begin{quote}
\ta{பழைய} \ta{வீட்டின்} \ta{மூலைகளில்} \ta{ஒட்டடை} \ta{நிறைந்திருந்தது}.
\end{quote}

\begin{quote}
\emph{paḻaiya vīṭṭiṉ mūlaikaḷil oṭṭaṭai niṟaintiruntatu.}
\end{quote}

\begin{quote}
The corners of the old house were full of cobwebs.
\end{quote}

To clear cobwebs is \ta{ஒட்டடை} \ta{அடி}, with a long-handled brush.

\begin{cousinweb}
The origin of \textbf{\ta{ஒட்டடை}} (\emph{oṭṭaṭai}) is not settled: no article of the Dravidian dictionary gives it, though it looks native.
\end{cousinweb}

\begin{grammarlens}[title={What you've built}]
One more word for describing what you lived and felt.
\end{grammarlens}

\subsection*{Your turn}
\begin{itemize}
\raggedright
  \item \emph{Say it:} \emph{oṭṭaṭai}
  \item \emph{Say it:} \emph{oṭṭaṭai}, once more
  \item \emph{Say it:} \emph{paḻaiya vīṭṭiṉ mūlaikaḷil oṭṭaṭai niṟaintiruntatu.} — The corners of the old house were full of cobwebs.
  \item \emph{Recall:} say the Tamil for a drawer, then the Tamil for a bedspread, then say \emph{oṭṭaṭai} again
\end{itemize}

\begin{tcolorbox}[breakable,colback=teal!4,colframe=teal!35!black,title={Before you move on}]
What does \ta{ஒட்டடை} mean? (\textquotedblleft{}A cobweb\textquotedblright{}.) Say it once more.
\end{tcolorbox}
\section[\texorpdfstring{\ta{தாழ்ப்பாள்} (tāḻppāḷ)}{தாழ்ப்பாள் (tāḻppāḷ)}]{\ta{தாழ்ப்பாள்} (tāḻppāḷ) — a latch; a door bolt}

\label{lesson:TA-C521-talppal}



\begin{quote}
Before the new one: say the Tamil for a bedspread, then the Tamil for a cobweb.
\end{quote}

\subsection*{You'll want to know: \ta{தாழ்ப்பாள்}}
\textbf{\ta{தாழ்ப்பாள்}} — \emph{tāḻppāḷ} — \textquotedblleft{}a latch; a door bolt\textquotedblright{}.

\begin{quote}
\ta{தூங்கும்} \ta{முன்} \ta{அப்பா} \ta{கதவின்} \ta{தாழ்ப்பாளைப்} \ta{போட்டார்}.
\end{quote}

\begin{quote}
\emph{tūṅkum muṉ appā kataviṉ tāḻppāḷaip pōṭṭār.}
\end{quote}

\begin{quote}
Before sleeping, Father bolted the door.
\end{quote}

One puts a bolt on with \ta{போடு} (\emph{pōṭu}): \ta{தாழ்ப்பாள்} \ta{போடு}, to bolt the door.

\begin{cousinweb}
\textbf{\ta{தாழ்ப்பாள்}} (\emph{tāḻppāḷ}) is native Dravidian: DEDR 3179 gives \ta{தாழ்} (\emph{tāḻ}), \textquotedblleft{}a bolt, a latch\textquotedblright{}, and the compound \ta{தாழ்ப்பாள்} for a door bolt.
\end{cousinweb}

\begin{grammarlens}[title={What you've built}]
One more word for describing what you lived and felt.
\end{grammarlens}

\subsection*{Your turn}
\begin{itemize}
\raggedright
  \item \emph{Say it:} \emph{tāḻppāḷ}
  \item \emph{Say it:} \emph{tāḻppāḷ}, once more
  \item \emph{Say it:} \emph{tūṅkum muṉ appā kataviṉ tāḻppāḷaip pōṭṭār.} — Before sleeping, Father bolted the door.
  \item \emph{Recall:} say the Tamil for a bedspread, then the Tamil for a cobweb, then say \emph{tāḻppāḷ} again
\end{itemize}

\begin{tcolorbox}[breakable,colback=teal!4,colframe=teal!35!black,title={Before you move on}]
What does \ta{தாழ்ப்பாள்} mean? (\textquotedblleft{}A latch; a door bolt\textquotedblright{}.) Say it once more.
\end{tcolorbox}
\section[\texorpdfstring{\ta{பூந்தொட்டி} (pūntoṭṭi)}{பூந்தொட்டி (pūntoṭṭi)}]{\ta{பூந்தொட்டி} (pūntoṭṭi) — a flowerpot}

\label{lesson:TA-C521-puntotti}



\begin{quote}
Before the new one: say the Tamil for a cobweb, then the Tamil for a latch.
\end{quote}

\subsection*{You'll want to know: \ta{பூந்தொட்டி}}
\textbf{\ta{பூந்தொட்டி}} — \emph{pūntoṭṭi} — \textquotedblleft{}a flowerpot\textquotedblright{}.

\begin{quote}
\ta{பால்கனியில்} \ta{மூன்று} \ta{பூந்தொட்டிகள்} \ta{இருக்கின்றன}.
\end{quote}

\begin{quote}
\emph{pālkaṉiyil mūṉṟu pūntoṭṭikaḷ irukkiṉṟaṉa.}
\end{quote}

\begin{quote}
There are three flowerpots on the balcony.
\end{quote}

Plural: \ta{பூந்தொட்டிகள்}, \emph{pūntoṭṭikaḷ}.

\ta{பூ} \textquotedblleft{}a flower\textquotedblright{} + \ta{தொட்டி} (\emph{toṭṭi}) \textquotedblleft{}a tub\textquotedblright{}: \ta{பூ} adds \ta{ந்} before \ta{தொட்டி}.

\begin{cousinweb}
\textbf{\ta{பூந்தொட்டி}} (\emph{pūntoṭṭi}) is a native compound: \ta{பூ} (\emph{pū}), \textquotedblleft{}a flower\textquotedblright{} (DEDR 4345), and \ta{தொட்டி} (\emph{toṭṭi}), \textquotedblleft{}a tub, a tank\textquotedblright{}.
\end{cousinweb}

\begin{grammarlens}[title={What you've built}]
One more word for describing what you lived and felt.
\end{grammarlens}

\subsection*{Your turn}
\begin{itemize}
\raggedright
  \item \emph{Say it:} \emph{pūntoṭṭi}
  \item \emph{Say it:} \emph{pūntoṭṭi}, once more
  \item \emph{Say it:} \emph{pālkaṉiyil mūṉṟu pūntoṭṭikaḷ irukkiṉṟaṉa.} — There are three flowerpots on the balcony.
  \item \emph{Recall:} say the Tamil for a cobweb, then the Tamil for a latch, then say \emph{pūntoṭṭi} again
\end{itemize}

\begin{tcolorbox}[breakable,colback=teal!4,colframe=teal!35!black,title={Before you move on}]
What does \ta{பூந்தொட்டி} mean? (\textquotedblleft{}A flowerpot\textquotedblright{}.) Say it once more.
\end{tcolorbox}
\section[(dialogue)]{Looking back: from Unions and Warehouses to Around the House}

\label{lesson:TA-R521-review-514-521}



\begin{quote}
Say the Tamil for a trade union and a strike.
\end{quote}

\subsection*{Your turn}
Say these aloud:

\begin{itemize}
\raggedright
  \item a trade union, a strike, a workshop, a warehouse, a recipient
  \item technology, software, hardware, a laptop, an e-book
  \item a web browser, a search engine, a blog, a video, a user
  \item a battery, to charge, a wire, to disconnect, to glow
  \item a computer program, a code, a backup copy, a robot, artificial
  \item a loudspeaker, a microphone, a telephone, a printer, electronic
  \item a satellite, automatic, a washing machine, an electrical appliance, to iron
  \item a drawer, a bedspread, a cobweb, a latch, a flowerpot
\end{itemize}

\begin{tcolorbox}[breakable,colback=teal!4,colframe=teal!35!black,title={Before you move on}]
A trade union? (\textbf{\ta{தொழிற்சங்கம்}}, \emph{toḻiṟcaṅkam}.) A strike? (\textbf{\ta{வேலைநிறுத்தம்}}, \emph{vēlainiṟuttam}.) A workshop? (\textbf{\ta{பணிமனை}}, \emph{paṇimaṉai}.) A warehouse? (\textbf{\ta{கிடங்கு}}, \emph{kiṭaṅku}.) A recipient? (\textbf{\ta{பெறுநர்}}, \emph{peṟunar}.) Technology? (\textbf{\ta{தொழில்நுட்பம்}}, \emph{toḻilnuṭpam}.) Software? (\textbf{\ta{மென்பொருள்}}, \emph{meṉporuḷ}.) Hardware? (\textbf{\ta{வன்பொருள்}}, \emph{vaṉporuḷ}.) A laptop? (\textbf{\ta{மடிக்கணினி}}, \emph{maṭikkaṇiṉi}.) An e-book? (\textbf{\ta{மின்னூல்}}, \emph{miṉṉūl}.) A web browser? (\textbf{\ta{உலாவி}}, \emph{ulāvi}.) A search engine? (\textbf{\ta{தேடுபொறி}}, \emph{tēṭupoṟi}.) A blog? (\textbf{\ta{வலைப்பதிவு}}, \emph{valaippativu}.) A video? (\textbf{\ta{காணொளி}}, \emph{kāṇoḷi}.) A user? (\textbf{\ta{பயனர்}}, \emph{payaṉar}.) A battery? (\textbf{\ta{மின்கலம்}}, \emph{miṉkalam}.) To charge? (\textbf{\ta{மின்னேற்று}}, \emph{miṉṉēṟṟu}.) A wire? (\textbf{\ta{கம்பி}}, \emph{kampi}.) To disconnect? (\textbf{\ta{துண்டி}}, \emph{tuṇṭi}.) To glow? (\textbf{\ta{ஒளிர்}}, \emph{oḷir}.) A computer program? (\textbf{\ta{நிரல்}}, \emph{niral}.) A code? (\textbf{\ta{குறியீடு}}, \emph{kuṟiyīṭu}.) A backup copy? (\textbf{\ta{காப்புப்பிரதி}}, \emph{kāppuppirati}.) A robot? (\textbf{\ta{ரோபோ}}, \emph{rōpō}.) Artificial? (\textbf{\ta{செயற்கை}}, \emph{ceyaṟkai}.) A loudspeaker? (\textbf{\ta{ஒலிபெருக்கி}}, \emph{oliperukki}.) A microphone? (\textbf{\ta{ஒலிவாங்கி}}, \emph{olivāṅki}.) A telephone? (\textbf{\ta{தொலைபேசி}}, \emph{tolaipēci}.) A printer? (\textbf{\ta{அச்சுப்பொறி}}, \emph{accuppoṟi}.) Electronic? (\textbf{\ta{மின்னணு}}, \emph{miṉṉaṇu}.) A satellite? (\textbf{\ta{செயற்கைக்கோள்}}, \emph{ceyaṟkaikkōḷ}.) Automatic? (\textbf{\ta{தானியங்கி}}, \emph{tāṉiyaṅki}.) A washing machine? (\textbf{\ta{சலவை} \ta{இயந்திரம்}}, \emph{calavai iyantiram}.) An electrical appliance? (\textbf{\ta{மின்சாதனம்}}, \emph{miṉcātaṉam}.) To iron? (\textbf{\ta{இஸ்திரி} \ta{போடு}}, \emph{istiri pōṭu}.) A drawer? (\textbf{\ta{இழுப்பறை}}, \emph{iḻuppaṟai}.) A bedspread? (\textbf{\ta{விரிப்பு}}, \emph{virippu}.) A cobweb? (\textbf{\ta{ஒட்டடை}}, \emph{oṭṭaṭai}.) A latch? (\textbf{\ta{தாழ்ப்பாள்}}, \emph{tāḻppāḷ}.) A flowerpot? (\textbf{\ta{பூந்தொட்டி}}, \emph{pūntoṭṭi}.)
\end{tcolorbox}
